% Source: https://artofproblemsolving.com/community/c3h1566372_solving_linear_equations
% Stable identity: aops-topic-1566372-q01
If two equations $P=Q$ and $R=T$ both hold, then $PR=QT$ holds, and so does $\dfrac {P}{R}=\dfrac {Q}{T}$ provided $R\ne0$: we are just multiplying or dividing equal numbers by equal numbers. The only care needed is that the new equation alone may not be equivalent to the system, so the solutions found must be checked in the original equations.

Here $y\ne0$ (it is a denominator), and $x\ne0$ because $xy=12$.

\textbf{Multiply:} $xy\cdot\dfrac{x}{y}=12\cdot3$, so $x^2=36$ and $x=\pm6$.

\textbf{Divide} the first equation by the second (allowed, as $\dfrac {x}{y}=3\ne0$): $xy\div\dfrac {x}{y}=y^2$, so $y^2=\dfrac{12}{3}=4$ and $y=\pm2$.

The second equation says $x=3y$, so $x$ and $y$ have the same sign. Checking the four sign combinations leaves
\[
(x,y)=\boxed{(6,2)\ \text{or}\ (-6,-2)},
\]
and both satisfy $xy=12$ and $\dfrac {x}{y}=3$. (The pairs $(6,-2)$ and $(-6,2)$ satisfy $x^2=36$ and $y^2=4$ but not the original system; this is why the check is needed.)
